% =====================================================================
\chapter{Limites}

O conceito de limite está na base de todo o Cálculo: a derivada e a integral são, ambas,
definidas por limites. Neste capítulo construímos a ideia de limite, primeiro de modo
intuitivo, com tabelas e gráficos, depois com regras de cálculo e, por fim, com uma
definição precisa. Também estudamos a continuidade, uma propriedade que depende
diretamente do limite.

% ---------------------------------------------------------------------
\section{Uma prévia do Cálculo}

O Cálculo nasceu de dois problemas geométricos: encontrar a reta tangente a uma curva
(o \emph{problema da tangente}, que leva ao Cálculo Diferencial) e calcular a área sob uma
curva (o \emph{problema da área}, que leva ao Cálculo Integral). Os dois envolvem uma
mesma ideia: aproximar uma quantidade cada vez melhor e observar para onde as aproximações
tendem.

\subsection{O problema da tangente}

Para uma reta, a taxa de variação é a inclinação. Para uma curva qualquer, aproximamos a
taxa de variação em um ponto $(a,f(a))$ por \emph{retas secantes}, que passam por
$(a,f(a))$ e por um ponto próximo $(x,f(x))$.

\begin{definicao}
A \textbf{reta secante}\index{reta!secante} ao gráfico de $f$ pelos pontos $(a,f(a))$ e $(x,f(x))$ é a reta
que passa por esses pontos. Sua inclinação é
\[
m_{\sec}=\frac{f(x)-f(a)}{x-a}.
\]
\end{definicao}

À medida que $x$ se aproxima de $a$, as secantes se aproximam de uma reta, a
\textbf{reta tangente}\index{reta!tangente} ao gráfico em $(a,f(a))$, e $m_{\sec}$ se aproxima da inclinação
da tangente, que é a \emph{taxa de variação instantânea} de $f$ em $a$.

\figtikz{\begin{tikzpicture}
\begin{axis}[painel,width=7cm,height=5.2cm,xmin=-0.3,xmax=3.6,ymin=0,ymax=3.6,xtick={1},xticklabels={$a$},ytick=\empty]
\addplot[domain=-0.3:3.4]{0.25*x^2+1} node[pos=0.98,left,font=\scriptsize]{$f$};
\foreach \b/\c in {3/laranja!90,2.2/laranja!60,1.5/laranja!35}{\edef\tmp{\noexpand\addplot[\c,thin,domain=-0.3:3.6]{1.25+0.25*(\b+1)*(x-1)};}\tmp}
\addplot[verde,thick,domain=-0.3:3.6]{1.25+0.5*(x-1)} node[pos=0.96,below,font=\scriptsize]{tangente};
\fill (axis cs:1,1.25) circle (1.6pt);
\end{axis}\end{tikzpicture}}{Quando $x$ se aproxima de $a$, as secantes (laranja) se aproximam da reta tangente (verde).}

\begin{exemplo}{2.1}{Inclinações de secantes}
Estime a inclinação da tangente a $f(x)=x^2$ em $x=1$ usando as secantes por $(1,1)$ e
(a) $(2,4)$; (b) $(3/2,\,9/4)$.
\solucao
(a) $m_{\sec}=\dfrac{4-1}{2-1}=3$. (b) $m_{\sec}=\dfrac{9/4-1}{3/2-1}=\dfrac52=2{,}5$. O
ponto de (b) está mais perto de $(1,1)$, então $2{,}5$ é uma estimativa melhor. Com pontos
cada vez mais próximos, veremos que a inclinação da tangente é exatamente 2.
\end{exemplo}

O mesmo raciocínio se aplica ao movimento. Se $s(t)$ é a posição de um objeto no instante
$t$, a \textbf{velocidade média}\index{velocidade!media@média} no intervalo $[a,t]$ é
\[
v_{\text{méd}}=\frac{s(t)-s(a)}{t-a},
\]
e a \textbf{velocidade instantânea}\index{velocidade!instantanea@instantânea} em $t=a$ é o valor do qual as velocidades médias se
aproximam quando $t$ se aproxima de $a$, se esse valor existir.

\subsection{O problema da área}

Para estimar a área sob o gráfico de uma função positiva entre $x=a$ e $x=b$, dividimos o
intervalo em partes e somamos as áreas de retângulos. Quanto mais retângulos, melhor a
aproximação; a área exata é o valor para o qual essas somas tendem. Esse é o ponto de
partida do Cálculo Integral (capítulo~5).

\figtikz{\begin{tikzpicture}
\begin{axis}[painel,width=7cm,height=5cm,xmin=0,xmax=5.6,ymin=0,ymax=2.6,xtick={1,5},xticklabels={$a$,$b$},ytick=\empty]
\foreach \a in {1,1.5,...,4.5}{\edef\tmp{\noexpand\draw[fill=azul!15,draw=azul] (axis cs:\a,0) rectangle (axis cs:{\a+0.5},{-0.15*(\a-3.2)^2+2.2});}\tmp}
\addplot[domain=0.2:5.5]{-0.15*(x-3.2)^2+2.2};
\end{axis}\end{tikzpicture}}{A área sob a curva aproximada por retângulos.}

% ---------------------------------------------------------------------
\section{O limite de uma função}

\subsection{Definição intuitiva}

\begin{definicao}
Seja $f$ definida em todos os pontos de um intervalo aberto contendo $a$, exceto
possivelmente em $a$, e seja $L$ um número real. Se os valores $f(x)$ se aproximam de $L$
quando $x$ ($\neq a$) se aproxima de $a$, dizemos que o \textbf{limite}\index{limite} de $f(x)$ quando
$x$ tende a $a$ é $L$, e escrevemos
\[
\lim_{x\to a}f(x)=L .
\]
Em outras palavras: à medida que $x$ fica perto de $a$, $f(x)$ fica perto de $L$ e
permanece perto.
\end{definicao}

\begin{nota}
O limite não depende do valor $f(a)$; $f$ nem precisa estar definida em $a$. Importa
apenas o comportamento de $f$ \emph{perto} de $a$.
\end{nota}

\begin{estrategia}{estimar um limite com uma tabela}\index{limite!estimado por tabela}
\begin{enumerate}[label=\arabic*.]
\item Monte duas colunas de valores de $x$ que se aproximam de $a$: uma pela esquerda
  ($a-0{,}1$, $a-0{,}01$, $a-0{,}001,\dots$) e outra pela direita ($a+0{,}1$,
  $a+0{,}01,\dots$), com os respectivos $f(x)$.
\item Se as duas colunas de $f(x)$ se aproximam de um mesmo valor $L$, a estimativa é
  $\lim_{x\to a}f(x)=L$.
\item Confirme com o gráfico de $f$ perto de $x=a$.
\end{enumerate}
\end{estrategia}

\begin{exemplo}{2.2}{Estimando $\lim_{x\to0}\frac{\sin x}{x}$}
Calculando $(\sin x)/x$ para $x=\pm0{,}1;\ \pm0{,}01;\ \pm0{,}001$, obtemos
$0{,}99833\ldots$, $0{,}99998\ldots$ e $0{,}9999998\ldots$ (a função é par, então os valores
pela esquerda e pela direita coincidem). Os valores se aproximam de 1, e o gráfico confirma:
\[
\lim_{x\to0}\frac{\sin x}{x}=1 .
\]
Provaremos esse resultado na seção~2.3.
\figtikz{\begin{tikzpicture}
\begin{axis}[painel,width=8cm,height=4.4cm,xmin=-8,xmax=8,ymin=-0.45,ymax=1.2,xtick={-6,-3,3,6},ytick={0.5,1}]
\addplot[domain=-8:-0.01,samples=150]{sin(deg(x))/x}; \addplot[domain=0.01:8,samples=150]{sin(deg(x))/x};
\draw[azul,fill=white] (axis cs:0,1) circle (2pt);
\end{axis}\end{tikzpicture}}{Gráfico de $f(x)=(\sin x)/x$ perto de $x=0$.}
\end{exemplo}

Dois limites básicos, intuitivamente evidentes, servem de ponto de partida:

\begin{teorema}[Dois limites importantes]{2.1}
Para todo número real $a$ e toda constante $c$:
\qquad $\displaystyle\lim_{x\to a}x=a$ \qquad e \qquad $\displaystyle\lim_{x\to a}c=c$.
\end{teorema}

\subsection{Quando o limite não existe}

Um limite pode não existir por diferentes motivos: a função pode se aproximar de valores
diferentes pela esquerda e pela direita, crescer sem limitação, ou oscilar.

\begin{exemplo}{2.3}{Um limite que não existe}
Considere $\lim_{x\to0}\sin(1/x)$. Para a sequência $x=\dfrac2\pi,\dfrac2{3\pi},
\dfrac2{5\pi},\dots$, que tende a 0, os valores de $\sin(1/x)$ são $1,-1,1,-1,\dots$. Eles não
se aproximam de nenhum número, logo o limite não existe. O gráfico oscila cada vez mais
rapidamente entre $-1$ e $1$ perto de $x=0$.
\figtikz{\begin{tikzpicture}
\begin{axis}[painel,width=8cm,height=4.4cm,xmin=-0.5,xmax=0.5,ymin=-1.3,ymax=1.3,xtick={-0.4,-0.2,0.2,0.4},ytick={-1,1},
  xticklabel style={/pgf/number format/use comma}]
\addplot[thin,domain=-0.5:-0.012,samples=1200]{sin(deg(1/x))}; \addplot[thin,domain=0.012:0.5,samples=1200]{sin(deg(1/x))};
\end{axis}\end{tikzpicture}}{$\sin(1/x)$ oscila entre $-1$ e $1$ quando
$x\to0$.}
\end{exemplo}

\subsection{Limites laterais}

\begin{definicao}
\textbf{Limite à esquerda:}\index{limite!lateral} se $f$ está definida em um intervalo $(c,a)$ e $f(x)$ se
aproxima de $L$ quando $x$ se aproxima de $a$ com $x<a$, escrevemos
$\displaystyle\lim_{x\to a^-}f(x)=L$.

\textbf{Limite à direita:} se $f$ está definida em um intervalo $(a,c)$ e $f(x)$ se
aproxima de $L$ quando $x$ se aproxima de $a$ com $x>a$, escrevemos
$\displaystyle\lim_{x\to a^+}f(x)=L$.
\end{definicao}

\begin{exemplo}{2.4}{Limites laterais}
Para $f(x)=\begin{cases}x+1, & x<2\\ x^2-4, & x\ge2,\end{cases}$ temos
$\lim_{x\to2^-}f(x)=2+1=3$ e $\lim_{x\to2^+}f(x)=2^2-4=0$. Como os limites laterais são
diferentes, $\lim_{x\to2}f(x)$ não existe.
\figtikz{\begin{tikzpicture}
\begin{axis}[painel,width=6.2cm,height=5cm,xmin=-3,xmax=4,ymin=-2.2,ymax=5.5,xtick={-1,2},ytick={3}]
\addplot[domain=-3:2]{x+1}; \addplot[laranja,domain=2:3.05]{x^2-4};
\draw[azul,fill=white] (axis cs:2,3) circle (2pt); \fill[laranja] (axis cs:2,0) circle (2pt);
\end{axis}\end{tikzpicture}}{Limites laterais diferentes em $x=2$.}
\end{exemplo}

\begin{teorema}[Limites laterais e bilaterais]{2.2}
Seja $f$ definida em um intervalo aberto contendo $a$, exceto possivelmente em $a$. Então
\[
\lim_{x\to a}f(x)=L \iff \lim_{x\to a^-}f(x)=L \ \text{ e } \ \lim_{x\to a^+}f(x)=L .
\]
\end{teorema}

\subsection{Limites infinitos}

Quando $f(x)$ cresce sem limitação à medida que $x$ se aproxima de $a$ (pela esquerda,
pela direita ou pelos dois lados), escrevemos, respectivamente,
\[
\lim_{x\to a^-}f(x)=+\infty,\qquad \lim_{x\to a^+}f(x)=+\infty,\qquad
\lim_{x\to a}f(x)=+\infty ;
\]
analogamente para $-\infty$, quando $f(x)$ decresce sem limitação. Atenção: dizer que o
limite é $\infty$ é uma forma de descrever \emph{como} o limite deixa de existir, pois
$\infty$ não é um número real. Por exemplo, para $f(x)=1/x$:
$\lim_{x\to0^-}\frac1x=-\infty$ e $\lim_{x\to0^+}\frac1x=+\infty$.

\figtikz{\begin{tikzpicture}
\begin{axis}[painel,width=6cm,height=5cm,xmin=-3,xmax=3,ymin=-5,ymax=5,xtick={-2,-1,1,2},ytick={-4,-2,2,4}]
\addplot[domain=-3:-0.2]{1/x}; \addplot[domain=0.2:3]{1/x};
\end{axis}\end{tikzpicture}}{O gráfico de $1/x$ perto de $x=0$.}

\begin{teorema}[Limites infinitos de potências inteiras]{2.3}\index{limite!infinito}
Se $n$ é um inteiro positivo par, $\displaystyle\lim_{x\to a}\frac{1}{(x-a)^n}=+\infty$.
Se $n$ é um inteiro positivo ímpar,
$\displaystyle\lim_{x\to a^+}\frac{1}{(x-a)^n}=+\infty$ e
$\displaystyle\lim_{x\to a^-}\frac{1}{(x-a)^n}=-\infty$.
\end{teorema}

\begin{definicao}
A reta $x=a$ é uma \textbf{assíntota vertical}\index{assintota@assíntota!vertical} de $f$ se pelo menos um dos limites
$\lim_{x\to a^-}f(x)$, $\lim_{x\to a^+}f(x)$ ou $\lim_{x\to a}f(x)$ for $+\infty$ ou
$-\infty$.
\end{definicao}

Por exemplo, $f(x)=\dfrac{1}{(x-3)^4}$ tem assíntota vertical $x=3$, pois
$\lim_{x\to3}f(x)=+\infty$ pelo Teorema~2.3.

% ---------------------------------------------------------------------
\section{As propriedades dos limites}

\subsection{Calculando limites com as propriedades}

Estimar limites com tabelas e gráficos é útil, mas trabalhoso e às vezes enganoso. As
propriedades a seguir permitem calcular limites diretamente.

\begin{teorema}[Propriedades dos limites]{2.4}
Sejam $f$ e $g$ definidas para $x\neq a$ em um intervalo aberto contendo $a$, com
$\lim_{x\to a}f(x)=L$ e $\lim_{x\to a}g(x)=M$ reais, e seja $c$ uma constante. Então:
\begin{itemize}
\item \textbf{Soma e diferença:} $\lim_{x\to a}\big(f(x)\pm g(x)\big)=L\pm M$.
\item \textbf{Múltiplo constante:} $\lim_{x\to a}c\,f(x)=cL$.
\item \textbf{Produto:} $\lim_{x\to a}f(x)\,g(x)=L\cdot M$.
\item \textbf{Quociente:} $\lim_{x\to a}\dfrac{f(x)}{g(x)}=\dfrac LM$, se $M\neq0$.
\item \textbf{Potência:} $\lim_{x\to a}\big(f(x)\big)^n=L^n$, para todo inteiro positivo $n$.
\item \textbf{Raiz:} $\lim_{x\to a}\sqrt[n]{f(x)}=\sqrt[n]{L}$, para todo $L$ se $n$ é ímpar,
  e para $L\ge0$ se $n$ é par e $f(x)\ge0$.
\end{itemize}
\end{teorema}

\begin{exemplo}{2.5}{Usando as propriedades}
\[
\lim_{x\to2}\frac{2x^2-3x+1}{x^3+4}
=\frac{2\lim x^2-3\lim x+\lim 1}{\lim x^3+\lim 4}
=\frac{2(4)-3(2)+1}{8+4}=\frac{3}{12}=\frac14 .
\]
\end{exemplo}

Aplicando as propriedades repetidamente, obtemos o resultado geral para polinômios e
funções racionais.

\begin{teorema}[Limites de polinômios e funções racionais]{2.5}
Se $p$ e $q$ são polinômios e $a$ é real, então
$\displaystyle\lim_{x\to a}p(x)=p(a)$ e
$\displaystyle\lim_{x\to a}\frac{p(x)}{q(x)}=\frac{p(a)}{q(a)}$ quando $q(a)\neq0$.
\end{teorema}

\subsection{A forma indeterminada $0/0$}

Quando $q(a)=0$, a propriedade do quociente não se aplica diretamente. Se também $p(a)=0$,
temos a \textbf{forma indeterminada $0/0$}: o limite pode existir ou não, e é preciso
manipular a expressão. A ideia central é que \emph{se $f(x)=g(x)$ para todo $x\neq a$
perto de $a$, então as duas funções têm o mesmo limite em $a$}. Por exemplo,
$\dfrac{x^2-1}{x-1}=x+1$ para $x\neq1$, logo
$\lim_{x\to1}\dfrac{x^2-1}{x-1}=\lim_{x\to1}(x+1)=2$.

\begin{estrategia}{limites da forma $0/0$}\index{forma indeterminada}
\begin{enumerate}[label=\arabic*.]
\item Verifique que as propriedades não se aplicam diretamente (forma $0/0$).
\item Encontre uma função igual a $f(x)/g(x)$ para $x\neq a$:
  (a) se forem polinômios, fatore e cancele os fatores comuns;
  (b) se houver diferença com raiz quadrada, multiplique pelo conjugado;
  (c) se houver fração composta, simplifique-a.
\item Aplique as propriedades dos limites.
\end{enumerate}
\end{estrategia}

\begin{exemplo}{2.6}{Fatorando e cancelando}
\[
\lim_{x\to3}\frac{x^2-3x}{2x^2-5x-3}=\lim_{x\to3}\frac{x(x-3)}{(x-3)(2x+1)}
=\lim_{x\to3}\frac{x}{2x+1}=\frac37 .
\]
\end{exemplo}

\begin{exemplo}{2.7}{Multiplicando pelo conjugado}
Calcule $\displaystyle\lim_{x\to-1}\frac{\sqrt{x+2}-1}{x+1}$.
\solucao
A expressão tem a forma $0/0$ em $x=-1$. Multiplicamos numerador e denominador pelo
conjugado $\sqrt{x+2}+1$:
\begin{align*}
\lim_{x\to-1}\frac{\sqrt{x+2}-1}{x+1}
&=\lim_{x\to-1}\frac{(\sqrt{x+2}-1)(\sqrt{x+2}+1)}{(x+1)(\sqrt{x+2}+1)}
 =\lim_{x\to-1}\frac{x+1}{(x+1)(\sqrt{x+2}+1)}\\
&=\lim_{x\to-1}\frac{1}{\sqrt{x+2}+1}=\frac12 .
\end{align*}
\end{exemplo}

\begin{exemplo}{2.8}{Simplificando uma fração composta}
\[
\lim_{x\to1}\frac{\frac{1}{x+1}-\frac12}{x-1}
=\lim_{x\to1}\frac{2-(x+1)}{2(x+1)(x-1)}
=\lim_{x\to1}\frac{-(x-1)}{2(x+1)(x-1)}
=\lim_{x\to1}\frac{-1}{2(x+1)}=-\frac14 .
\]
\end{exemplo}

\textbf{Forma $K/0$ com $K\neq0$.} Se o numerador tende a $K\neq0$ e o denominador tende a
0, o limite não é finito. Estudamos os limites laterais pelo sinal: por exemplo, em
$\lim_{x\to2}\frac{x-3}{x^2-2x}$, o numerador tende a $-1$ e o denominador $x(x-2)$ tende
a $0^-$ pela esquerda e a $0^+$ pela direita. Logo o limite é $+\infty$ pela esquerda e
$-\infty$ pela direita.

\subsection{O Teorema do Confronto}

\begin{teorema}[Teorema do Confronto]{2.6}
Sejam $f$, $g$ e $h$ definidas para $x\neq a$ em um intervalo aberto contendo $a$. Se
$f(x)\le g(x)\le h(x)$ nesse intervalo, para $x\neq a$, e
$\lim_{x\to a}f(x)=L=\lim_{x\to a}h(x)$, então $\lim_{x\to a}g(x)=L$.
\end{teorema}

\figtikz{\begin{tikzpicture}
\begin{axis}[painel,width=7cm,height=5cm,xmin=-1,xmax=3.2,ymin=-0.5,ymax=3.2,xtick={1},xticklabels={$a$},ytick=\empty]
\addplot[domain=-1:3]{2+0.5*(x-1)^2} node[pos=0.97,above left,font=\scriptsize]{$h$};
\addplot[verde,domain=-1:3]{2-0.3*(x-1)^2+0.15*(x-1)^2*sin(deg(4*x))} node[pos=0.97,left,font=\scriptsize]{$g$};
\addplot[laranja,domain=-0.6:2.6]{2-(x-1)^2} node[pos=0.97,left,font=\scriptsize]{$f$};
\fill (axis cs:1,2) circle (1.6pt);
\end{axis}\end{tikzpicture}}{O Teorema do Confronto: $g$ fica ``espremida''
entre $f$ e $h$.}

\begin{exemplo}{2.9}{Aplicando o Teorema do Confronto}
Como $-1\le\cos x\le1$, temos $-|x|\le x\cos x\le|x|$. Como $\lim_{x\to0}(-|x|)=0=
\lim_{x\to0}|x|$, concluímos que $\lim_{x\to0}x\cos x=0$.
\end{exemplo}

\textbf{Dois limites trigonométricos fundamentais.} Pelo círculo unitário, para
$0<\theta<\pi/2$, valem $0<\sin\theta<\theta<\tan\theta$ (o arco de comprimento $\theta$
fica entre o segmento vertical de comprimento $\sin\theta$ e o segmento tangente de
comprimento $\tan\theta$). Pelo confronto, isso dá $\lim_{\theta\to0}\sin\theta=0$ e,
daí, $\lim_{\theta\to0}\cos\theta=1$.

\figtikz{\begin{tikzpicture}[scale=2.3,font=\small]
\draw[->] (-0.2,0) -- (1.35,0) node[right]{$x$}; \draw[->] (0,-0.2) -- (0,1.15) node[above]{$y$};
\draw[thick] (0,0) ++(0:1) arc (0:95:1);
\fill[azul!15] (0,0) -- (1,0) arc (0:40:1) -- cycle;
\draw[thick,azul] (0,0) -- (40:1); \draw[laranja,thick] (40:1) -- ({cos(40)},0); \draw[laranja,thin] ({cos(40)},0.32) -- (0.45,0.62) node[above,font=\footnotesize]{$\sin\theta$};
\draw[thick] (0,0) -- (1,{tan(40)}); \draw[verde,thick] (1,0) -- (1,{tan(40)}) node[midway,right]{$\tan\theta$};
\draw[azul,very thick] (1,0) arc (0:40:1); \draw (0.22,0) arc (0:40:0.22); \node at (20:0.33) {$\theta$};
\node[below] at (1,0) {$1$};
\end{tikzpicture}}{Para $0<\theta<\pi/2$: $\sin\theta<\theta<\tan\theta$ (segmento laranja, arco azul e segmento verde).}

Dividindo $\sin\theta<\theta<\tan\theta$ por $\sin\theta>0$ e invertendo, obtemos
$\cos\theta<\dfrac{\sin\theta}{\theta}<1$. Essa desigualdade vale também para
$-\pi/2<\theta<0$, porque as três funções são pares. Pelo Teorema do Confronto:

\begin{teorema}[Limites trigonométricos especiais]{2.7}
\[
\lim_{\theta\to0}\frac{\sin\theta}{\theta}=1
\qquad\text{e}\qquad
\lim_{\theta\to0}\frac{1-\cos\theta}{\theta}=0 .
\]
\end{teorema}
O segundo limite sai do primeiro multiplicando pelo conjugado:
\[
\frac{1-\cos\theta}{\theta}=\frac{1-\cos^2\theta}{\theta(1+\cos\theta)}
=\frac{\sin\theta}{\theta}\cdot\frac{\sin\theta}{1+\cos\theta}\ \longrightarrow\ 1\cdot\frac02=0 .
\]
Esses dois limites serão essenciais para derivar seno e cosseno no capítulo~3.

% ---------------------------------------------------------------------
\section{Continuidade}\index{continuidade}

\subsection{Continuidade em um ponto}

Intuitivamente, uma função é contínua quando seu gráfico pode ser traçado sem tirar o lápis
do papel. A definição precisa exige três condições.

\begin{definicao}
Uma função $f$ é \textbf{contínua no ponto $a$} se, e somente se:
\begin{enumerate}[label=(\roman*)]
\item $f(a)$ está definida;
\item $\lim_{x\to a}f(x)$ existe;
\item $\lim_{x\to a}f(x)=f(a)$.
\end{enumerate}
Se alguma condição falha, $f$ é \textbf{descontínua} em $a$.
\end{definicao}

Para verificar a continuidade em $a$, testamos as condições na ordem: se uma falha,
paramos. Por exemplo, $f(x)=(x^2-4)/(x-2)$ não é contínua em $x=2$, pois $f(2)$ não está
definida. Já
\[
f(x)=\begin{cases}x^3, & x\neq1\\ 0, & x=1\end{cases}
\]
está definida em 1 e tem $\lim_{x\to1}f(x)=1$, mas $1\neq f(1)=0$: falha a condição (iii).

\begin{teorema}[Continuidade de polinômios e funções racionais]{2.8}
Polinômios e funções racionais são contínuos em todos os pontos de seus domínios.
\end{teorema}
Isso decorre do Teorema~2.5: para $q(a)\neq0$, $\lim_{x\to a}p(x)/q(x)=p(a)/q(a)$.

\subsection{Tipos de descontinuidade}

\begin{definicao}
Se $f$ é descontínua em $a$, então $f$ tem:
\begin{enumerate}[label=\arabic*.]
\item uma \textbf{descontinuidade removível}\index{descontinuidade} se $\lim_{x\to a}f(x)$ existe (é um número
  real);
\item uma \textbf{descontinuidade de salto} se os limites laterais existem, mas são
  diferentes;
\item uma \textbf{descontinuidade infinita} se pelo menos um limite lateral é $\pm\infty$.
\end{enumerate}
\end{definicao}

\figtikz{\begin{tikzpicture}
\begin{groupplot}[group style={group size=3 by 1,horizontal sep=1.1cm},painel,width=5.2cm,height=4.4cm,
  xmin=-0.5,xmax=4,ymin=-0.5,ymax=4,xtick={2},xticklabels={$a$},ytick=\empty]
\nextgroupplot[title={(a) removível}]
\addplot[domain=-0.5:4]{0.5*x+1};
\draw[azul,thick,fill=white] (axis cs:2,2) circle (2pt);
\fill[laranja] (axis cs:2,3.2) circle (2pt);
\nextgroupplot[title={(b) de salto}]
\addplot[domain=-0.5:2]{0.5*x+0.6};
\addplot[domain=2:4]{0.5*x+1.6};
\fill[azul] (axis cs:2,1.6) circle (2pt);
\draw[azul,thick,fill=white] (axis cs:2,2.6) circle (2pt);
\nextgroupplot[title={(c) infinita},ymin=-4,ymax=4]
\addplot[domain=-0.5:1.75]{1/(x-2)+0.5};
\addplot[domain=2.25:4]{1/(x-2)+0.5};
\draw[assint] (axis cs:2,-4) -- (axis cs:2,4);
\end{groupplot}\end{tikzpicture}}{Tipos de descontinuidade em $x=a$: (a) removível, com
o limite existindo, mas diferente de $f(a)$ ou com $f(a)$ indefinido; (b) de salto, com
limites laterais finitos e diferentes; (c) infinita, com limite lateral $\pm\infty$.}

\begin{exemplo}{2.10}{Classificando descontinuidades}
(a) $f(x)=\dfrac{x^2-4}{x-2}$ em $x=2$: $\lim_{x\to2}f(x)=\lim(x+2)=4$ existe;
descontinuidade \emph{removível}.

(b) $f(x)=\begin{cases}-x^2+4, & x\le3\\ 4x-8, & x>3\end{cases}$ em $x=3$: os limites
laterais são $-5$ e $4$; descontinuidade \emph{de salto}.

(c) $f(x)=\dfrac{x+2}{x+1}$ em $x=-1$: pela esquerda o limite é $-\infty$ e pela direita
$+\infty$; descontinuidade \emph{infinita}.
\end{exemplo}

\subsection{Continuidade em um intervalo}

\begin{definicao}[Continuidade lateral]
$f$ é \textbf{contínua à direita}\index{continuidade!lateral} em $a$ se $\lim_{x\to a^+}f(x)=f(a)$, e
\textbf{contínua à esquerda} em $a$ se $\lim_{x\to a^-}f(x)=f(a)$.
\end{definicao}

Uma função é contínua em um intervalo aberto se é contínua em cada ponto dele. É contínua
em $[a,b]$ se é contínua em $(a,b)$, contínua à direita em $a$ e contínua à esquerda em
$b$. Por exemplo, $f(x)=\sqrt{4-x^2}$ é contínua em $[-2,2]$.

\begin{teorema}[Função composta]{2.9}
Se $f$ é contínua em $L$ e $\lim_{x\to a}g(x)=L$, então
\[
\lim_{x\to a}f\big(g(x)\big)=f\Big(\lim_{x\to a}g(x)\Big)=f(L).
\]
\end{teorema}
Por exemplo, $\lim_{x\to\pi/2}\cos\!\big(x-\tfrac\pi2\big)=\cos0=1$. Em particular, a
composta de funções contínuas é contínua.

\begin{teorema}[Continuidade das funções trigonométricas]{2.10}
As funções trigonométricas são contínuas em todo o seu domínio.
\end{teorema}

Mais geralmente, exponenciais, logaritmos, raízes e todas as funções construídas a partir
delas e de polinômios por operações algébricas e composição são contínuas onde estão
definidas.

\subsection{O Teorema do Valor Intermediário}

\begin{teorema}[Teorema do Valor Intermediário]{2.11}
Seja $f$ contínua em um intervalo fechado e limitado $[a,b]$. Se $z$ é qualquer número
real entre $f(a)$ e $f(b)$, então existe $c\in[a,b]$ tal que $f(c)=z$.
\end{teorema}

\figtikz{\begin{tikzpicture}
\begin{axis}[painel,width=7cm,height=5.4cm,xmin=0,xmax=5.6,ymin=0,ymax=5,xtick={0.3,4.61,5},xticklabels={$a$,$c$,$b$},ytick={1.3645,3,4.333},yticklabels={$f(a)$,$z$,$f(b)$}]
\draw[dashed] (axis cs:0,1.3645) -| (axis cs:0.3,0); \draw[dashed] (axis cs:0,4.333) -| (axis cs:5,0);
\draw[laranja,dashed] (axis cs:0,3) -| (axis cs:4.61,0);
\addplot[domain=0.3:5]{x^3/6-x^2+1.5*x+1};
\fill[laranja] (axis cs:4.61,3) circle (1.6pt);
\end{axis}\end{tikzpicture}}{Existe $c\in[a,b]$ com $f(c)=z$.}

\begin{exemplo}{2.11}{Garantindo um zero}
Mostre que $f(x)=x-\cos x$ tem pelo menos um zero.
\solucao
$f$ é contínua em $\R$. Como $f(0)=-1<0$ e $f(\pi/2)=\pi/2>0$, o Teorema do Valor
Intermediário garante que existe $c\in[0,\pi/2]$ com $f(c)=0$.
\end{exemplo}

\begin{nota}
A continuidade é essencial. Para $f(x)=1/x$, temos $f(-1)=-1<0$ e $f(1)=1>0$, mas $f$ não
tem zero em $[-1,1]$, pois não é contínua em $x=0$.
\end{nota}

% ---------------------------------------------------------------------
\section{A definição precisa de limite}\index{limite!definicao precisa (-)@definição precisa ($\varepsilon$-$\delta$)}

\subsection{Quantificando a proximidade}

A definição intuitiva fala em $f(x)$ ``se aproximar'' de $L$. Para torná-la precisa,
medimos distâncias: $|f(x)-L|$ é a distância entre $f(x)$ e $L$, e $|x-a|$ é a distância
entre $x$ e $a$. Dizer que $f(x)$ fica tão perto de $L$ quanto quisermos significa: para
qualquer tolerância $\varepsilon>0$, conseguimos garantir $|f(x)-L|<\varepsilon$ tomando
$x$ suficientemente perto de $a$, a uma distância menor que algum $\delta>0$.

\begin{definicao}[Definição $\varepsilon$-$\delta$ de limite]
Seja $f$ definida para $x\neq a$ em um intervalo aberto contendo $a$, e seja $L$ real.
Então $\lim_{x\to a}f(x)=L$ se, para todo $\varepsilon>0$, existe $\delta>0$ tal que
\[
0<|x-a|<\delta\ \Longrightarrow\ |f(x)-L|<\varepsilon .
\]
\end{definicao}

\figtikz{\begin{tikzpicture}
\begin{axis}[painel,width=8cm,height=6.2cm,xmin=0,xmax=4.2,ymin=0,ymax=4.6,
  xtick={1.6,2,2.4},xticklabels={$a-\delta$,$a$,$a+\delta$},
  ytick={1.95,2.6,3.25},yticklabels={$L-\varepsilon$,$L$,$L+\varepsilon$},
  x tick label style={font=\scriptsize,rotate=0},axis on top]
\fill[laranja!12] (axis cs:0,1.95) rectangle (axis cs:4.2,3.25);
\fill[azul!12] (axis cs:1.6,0) rectangle (axis cs:2.4,4.6);
\draw[laranja,dashed,thin] (axis cs:0,1.95) -- (axis cs:4.2,1.95);
\draw[laranja,dashed,thin] (axis cs:0,3.25) -- (axis cs:4.2,3.25);
\draw[azul!60,dashed,thin] (axis cs:1.6,0) -- (axis cs:1.6,4.6);
\draw[azul!60,dashed,thin] (axis cs:2.4,0) -- (axis cs:2.4,4.6);
\addplot[domain=0.2:4]{0.15*x^3-0.9*x^2+2.1*x+0.6};
\draw[azul,thick,fill=white] (axis cs:2,2.6) circle (2pt);
\end{axis}\end{tikzpicture}}{A definição $\varepsilon$-$\delta$: para cada faixa
horizontal $(L-\varepsilon,L+\varepsilon)$ existe uma faixa vertical $(a-\delta,a+\delta)$
na qual o gráfico, exceto possivelmente em $x=a$, fica dentro da faixa horizontal.}

\begin{estrategia}{provar que $\lim_{x\to a}f(x)=L$}
\begin{enumerate}[label=\arabic*.]
\item Comece com: ``Seja $\varepsilon>0$.''
\item Encontre $\delta$ (geralmente trabalhando ``de trás para a frente'' a partir de
  $|f(x)-L|<\varepsilon$) e escreva: ``Escolha $\delta=\ldots$''.
\item Escreva: ``Suponha $0<|x-a|<\delta$.''
\item Usando essa hipótese, mostre que $|f(x)-L|<\varepsilon$.
\item Conclua: ``Portanto, $\lim_{x\to a}f(x)=L$.''
\end{enumerate}
\end{estrategia}

\begin{exemplo}{2.12}{Uma prova $\varepsilon$-$\delta$ com função afim}
Prove que $\lim_{x\to-1}(4x+1)=-3$.
\solucao
Rascunho: queremos $|(4x+1)-(-3)|=|4x+4|=4|x+1|<\varepsilon$, ou seja, $|x+1|<\varepsilon/4$.
Isso sugere $\delta=\varepsilon/4$.

Prova: seja $\varepsilon>0$. Escolha $\delta=\varepsilon/4$. Suponha $0<|x+1|<\delta$.
Então $|(4x+1)-(-3)|=4|x+1|<4\delta=\varepsilon$. Portanto,
$\lim_{x\to-1}(4x+1)=-3$.
\end{exemplo}

\begin{exemplo}{2.13}{Uma prova $\varepsilon$-$\delta$ com função quadrática}
Prove que $\lim_{x\to-1}(x^2-2x+3)=6$.
\solucao
Rascunho: $|(x^2-2x+3)-6|=|x+1|\cdot|x-3|$. O fator $|x-3|$ depende de $x$, então
primeiro o limitamos: se $|x+1|<1$, então $-1<x+1<1$, isto é, $-5<x-3<-3$, e logo
$|x-3|<5$. Basta então ter também $|x+1|<\varepsilon/5$.

Prova: seja $\varepsilon>0$. Escolha $\delta=\min\{1,\varepsilon/5\}$. Suponha
$0<|x+1|<\delta$. Então $|x+1|<1$, donde $|x-3|<5$, e $|x+1|<\varepsilon/5$. Logo
$|(x^2-2x+3)-6|=|x+1|\,|x-3|<\tfrac{\varepsilon}{5}\cdot5=\varepsilon$.
\end{exemplo}

\subsection{Provando as propriedades dos limites}

A definição precisa permite demonstrar as propriedades da seção~2.3. Como exemplo,
provemos a da soma. Uma ferramenta recorrente é a \textbf{desigualdade triangular}\index{desigualdade triangular}:
$|a+b|\le|a|+|b|$ para quaisquer reais $a$ e $b$.

\emph{Prova da propriedade da soma.} Suponha $\lim_{x\to a}f(x)=L$ e
$\lim_{x\to a}g(x)=M$. Seja $\varepsilon>0$. Existe $\delta_1>0$ tal que
$0<|x-a|<\delta_1$ implica $|f(x)-L|<\varepsilon/2$, e existe $\delta_2>0$ tal que
$0<|x-a|<\delta_2$ implica $|g(x)-M|<\varepsilon/2$. Escolha $\delta=\min\{\delta_1,\delta_2\}$.
Se $0<|x-a|<\delta$, então
\[
|(f(x)+g(x))-(L+M)|\le|f(x)-L|+|g(x)-M|<\frac\varepsilon2+\frac\varepsilon2=\varepsilon. \qquad\square
\]

Para mostrar que um limite \emph{não} existe, basta exibir um $\varepsilon>0$ para o qual
nenhum $\delta$ funciona. Por exemplo, se $f(x)=-1$ para $x<0$ e $f(x)=1$ para $x\ge0$, com
$\varepsilon=1/2$ não há $\delta$ que sirva para nenhum candidato $L$: perto de 0 há valores
de $f$ iguais a $-1$ e a $1$, que distam 2 entre si.

\subsection{Limites laterais e infinitos}

As outras formas de limite têm definições análogas:
\begin{itemize}
\item $\lim_{x\to a^+}f(x)=L$ se, para todo $\varepsilon>0$, existe $\delta>0$ tal que
  $0<x-a<\delta\Rightarrow|f(x)-L|<\varepsilon$;
\item $\lim_{x\to a^-}f(x)=L$ se, para todo $\varepsilon>0$, existe $\delta>0$ tal que
  $-\delta<x-a<0\Rightarrow|f(x)-L|<\varepsilon$;
\item $\lim_{x\to a}f(x)=+\infty$ se, para todo $M>0$, existe $\delta>0$ tal que
  $0<|x-a|<\delta\Rightarrow f(x)>M$ (e $-\infty$ com $f(x)<-M$).
\end{itemize}

% ---------------------------------------------------------------------
\section*{Resumo do capítulo}\markright{Resumo do capítulo}
\addcontentsline{toc}{section}{Resumo do capítulo}
\begin{itemize}
\item A inclinação da tangente e a velocidade instantânea são limites de inclinações de
  secantes e de velocidades médias; a área sob uma curva é o limite de somas de áreas de
  retângulos.
\item $\lim_{x\to a}f(x)=L$ quando $f(x)$ fica arbitrariamente perto de $L$ para $x$
  suficientemente perto de $a$ ($x\neq a$); o valor $f(a)$ é irrelevante.
\item O limite bilateral existe se, e somente se, os laterais existem e são iguais.
\item Um limite infinito descreve uma assíntota vertical; $\infty$ não é um número.
\item Para polinômios e racionais (com denominador não nulo), basta substituir.
  Na forma $0/0$, fatore, use o conjugado ou simplifique antes.
\item Teorema do Confronto; $\lim_{\theta\to0}\frac{\sin\theta}{\theta}=1$ e
  $\lim_{\theta\to0}\frac{1-\cos\theta}{\theta}=0$.
\item Contínua em $a$: $f(a)$ definida, o limite existe e é igual a $f(a)$.
  Descontinuidades: removível, de salto, infinita.
\item Teorema do Valor Intermediário: uma função contínua em $[a,b]$ assume todos os
  valores entre $f(a)$ e $f(b)$.
\item Definição precisa: para todo $\varepsilon>0$ existe $\delta>0$ tal que
  $0<|x-a|<\delta\Rightarrow|f(x)-L|<\varepsilon$.
\end{itemize}

% Gerado por ferramentas/gerar.py; não edite à mão.
\section*{Exercícios selecionados}
\addcontentsline{toc}{section}{Exercícios selecionados}
Problemas escolhidos para desenvolver a resolução de problemas. As soluções completas,
com animações dos gráficos, estão no repositório do livro: \url{https://github.com/AleffAzevedo/calculo-1-pt/tree/main/exercicios/cap2}.

\begin{enumerate}[label=\textbf{\thechapter.\arabic*}, leftmargin=2.6em, itemsep=6pt]
\item \textbf{Da secante à tangente.} Calcule $\displaystyle\lim_{h\to0}\frac{\frac{1}{a+h}-\frac1a}{h}$, em que $a\neq0$ é constante.
Interprete o resultado em termos de retas secantes ao gráfico de $f(x)=1/x$.
\item \textbf{O limite de $\sin(ax)/x$.} (a) Use tabelas com $x=\pm0{,}1;\ \pm0{,}01;\ \pm0{,}001$ para estimar
$\lim_{x\to0}\frac{\sin2x}{x}$ e $\lim_{x\to0}\frac{\sin3x}{x}$.
(b) Conjecture o valor de $\lim_{x\to0}\frac{\sin ax}{x}$ para $a>0$ e prove-o com o
Teorema~2.7.
\item \textbf{Conjugado com raiz.} Mostre que a substituição direta leva a uma forma indeterminada e calcule
$\displaystyle\lim_{t\to9}\frac{t-9}{\sqrt t-3}$.
\item \textbf{Fatorando cubos.} Calcule $\displaystyle\lim_{x\to1}\frac{x^3-1}{x^2-1}$.
\item \textbf{Limites laterais infinitos.} Calcule $\displaystyle\lim_{x\to1^-}\frac{2x^2+7x-4}{x^2+x-2}$ e
$\displaystyle\lim_{x\to1^+}\frac{2x^2+7x-4}{x^2+x-2}$.
\item \textbf{Teorema do Confronto.} Calcule $\displaystyle\lim_{\theta\to0}\theta^2\cos\!\left(\frac1\theta\right)$.
\item \textbf{Escolhendo $k$ para a continuidade (I).} Encontre $k$ para que $f(x)=\begin{cases}3x+2, & x<k\\ 2x-3, & k\le x\le8\end{cases}$
seja contínua em todo o seu domínio.
\item \textbf{Escolhendo $k$ para a continuidade (II).} Encontre $k$ para que $f(x)=\begin{cases}e^{kx}, & 0\le x<4\\ x+3, & 4\le x\le8\end{cases}$
seja contínua em $[0,8]$.
\item \textbf{TVI: duas partículas.} Uma partícula se move numa reta com posição $s(t)$ contínua, e $s(2)=5$, $s(5)=2$. Outra
partícula tem posição $h(t)=s(t)-t$. Explique por que existe um instante $t\in[2,5]$ em que
$h(t)=0$. O que isso significa para a primeira partícula?
\item \textbf{TVI: localizando uma raiz.} Use o TVI para decidir se $2^x=x^3$ tem solução em $[1{,}25;\,1{,}375]$ ou em
$[1{,}375;\,1{,}5]$. Justifique para cada intervalo.
\item \textbf{Prova $\varepsilon$-$\delta$.} Use a definição precisa de limite para provar que $\displaystyle\lim_{x\to2}(x^2+2x)=8$.
\item \textbf{Provando que um limite não existe.} Use a definição precisa para provar que $\displaystyle\lim_{x\to1}\frac{|x-1|}{x-1}$ não existe.
\end{enumerate}

\subsection*{Respostas}
{\small\begin{enumerate}[label=\textbf{\thechapter.\arabic*}, leftmargin=2.6em, itemsep=2pt]
\item $-\dfrac{1}{a^2}$: é a inclinação da reta tangente a $y=1/x$ em $x=a$.
\item (a) 2 e 3; (b) $\lim_{x\to0}\frac{\sin ax}{x}=a$.
\item $0/0$; o limite é 6.
\item $\dfrac32$.
\item $-\infty$ pela esquerda e $+\infty$ pela direita.
\item 0, pelo Teorema do Confronto.
\item $k=-5$.
\item $k=\dfrac{\ln7}{4}\approx0{,}4865$.
\item $h(2)=3>0$ e $h(5)=-3<0$; pelo TVI, $h(c)=0$ para algum $c\in[2,5]$, ou seja, $s(c)=c$.
\item Com $g(x)=2^x-x^3$: $g(1{,}25)>0$ e $g(1{,}375)<0$, logo há raiz no 1º intervalo; no 2º, $g$ não muda de sinal e o TVI nada garante.
\item Escolha $\delta=\min\{1,\varepsilon/7\}$.
\item Tome $\varepsilon=\tfrac12$: perto de 1 a função vale $-1$ e $1$, que não cabem juntos em nenhuma faixa de largura 1.
\end{enumerate}}

